\documentclass[11pt]{article}

% ---------------------------------------------------------------------------
% Earned Autonomy --- an engineering-leader whitepaper on the 2026-07-28 MCP
% spec. Clean article class (no venue branding). Root owns preamble + metadata;
% one \input per section. Build: cd paper && latexmk -pdf main.tex
% ---------------------------------------------------------------------------

\usepackage[T1]{fontenc}
\usepackage{lmodern}
\usepackage[margin=1in]{geometry}
\usepackage{microtype}
\usepackage[numbers,sort&compress]{natbib}
\usepackage[hyphens]{url}
\usepackage{listings}
\usepackage{xcolor}
\usepackage{enumitem}
\usepackage{titlesec}
\usepackage[colorlinks=true,linkcolor=black,citecolor=black,urlcolor=blue]{hyperref}

% tighten section spacing to keep the page budget
\titlespacing*{\section}{0pt}{1.4ex plus 0.3ex minus 0.2ex}{0.9ex}

% listings style for the illustrative policy exhibit
\lstdefinestyle{policy}{%
  basicstyle=\ttfamily\small,
  frame=single,
  framerule=0.4pt,
  breaklines=true,
  showstringspaces=false,
  columns=fullflexible,
  keepspaces=true,
  xleftmargin=0.8em,
  xrightmargin=0.8em,
  aboveskip=0.6em,
  belowskip=0.4em,
}

\renewcommand{\abstractname}{Executive summary}

\title{\vspace{-1.2em}\textbf{Earned Autonomy}\\[0.35em]
  \large Conditioning Agent Trust on Delivered Outcomes, Not Task Accuracy}
\author{Andrew Zigler}
\date{July 2026}

\begin{document}
\maketitle

\begin{abstract}
\noindent Enterprises have adopted autonomous software agents faster than they
have learned to tell whether the agents are delivering. The dominant answer to
that trust problem, published under names like earned, graduated, and
progressive autonomy, conditions an agent's authority on how accurately it
performs its tasks. This paper argues that task accuracy is the wrong adaptive
signal. Autonomy conditioned on task accuracy re-creates at the trust layer the
same clean-record trap the audit layer already has: a fully authorized,
correctly typed, cleanly traced call can open a defective change while every
record reads clean, and an accuracy-graded ladder promotes a workflow at the
moment its output draws the least scrutiny. Delivered outcome, whether an
agent's change held in production, is the adaptive conditioning variable these
autonomy ladders leave out. It complements the preventive gates that screen
every change regardless of a workflow's record; it does not replace them. The
Model Context Protocol revision finalizing on 2026-07-28 is the first open
protocol standard to make per-action agent-tool authorization a first-class,
session-independent property, so a workflow's earned tier can be enforced per
action while its outcomes are evaluated on a cadence. The paper walks what the
release makes legible and enforceable, defines a \emph{delivery evidence layer}
that joins governed actions to their delivered outcomes, and describes earned
autonomy: a rolling trust cadence, modeled on the standard's own governance by
measurement, that grants and revokes agent authority on evidence.
\end{abstract}

\section{The question that predates the spec}

Organizations have poured budget, headcount, and roadmap into autonomous
software work, and most of them still cannot say whether it paid off. MIT Media
Lab's Project NANDA, surveying roughly three hundred enterprise generative-AI
initiatives, found that about 95\% showed no measurable P\&L return
\citep{nanda2025}. That figure has many causes. One of them is plain: a return
you cannot attribute is a return you cannot claim. The 95\% is, in part, an
attribution failure.

Adoption is not the constraint. Stanford's AI Index put enterprise AI adoption
at 78\% in its 2025 edition and 88\% a year later
\citep{stanford2025, stanford2026}. The tools are in the building. What remains
unanswered is whether the autonomous work they produce is delivering, and for
whom.

Scale is what makes the gap urgent. A single developer reviewing a single
assistant's suggestions is a manageable trust problem. A fleet of agents opening
changes across dozens of repositories, at a rate no human reads line by line, is
not. Trust at that scale has to be granted by policy rather than by inspection,
and policy needs a signal to run on. The governance question here is older than
any protocol. Before a wire format made agent activity legible, leaders already
needed to know which autonomous work had earned more scope and which had not.
The Model Context Protocol revision finalizing on 2026-07-28 does not create that
question. It changes what an organization can answer it with.

The question also already has a named answer in the field, and that answer is
worth taking seriously before improving on it. Earned autonomy, graduated
autonomy, and progressive autonomy all describe the same ladder: an agent starts
under human review, and as it performs, it earns the right to act with less
oversight. Matthias Roder's earned-autonomy gradient formalizes the rungs
\citep{roder}; a growing set of platform vendors ship progressive-autonomy
controls \citep{progressiveautonomy}; the Digital Apprentice line of work models
the same apprenticeship arc \citep{digitalapprentice}; and McKinsey frames agent
trust as decision rights granted and revoked over time \citep{mckinsey}. This
ladder is published across vendors and research, and
it is converging on a shared shape.

Every approach we surveyed conditions autonomy on the same class of signal: how
well the agent performed its task. Extraction accuracy, override rate, tool-call
correctness. That choice of variable is the one thing this paper revisits, because
it is the same variable that lets a clean record hide a bad outcome. This paper
keeps the name earned autonomy and changes only the variable it conditions on:
delivered outcome, not task accuracy. The rest of the argument shows why the
2026-07-28 spec is the first open protocol standard on which that condition can
be enforced per action.

\section{What the spec makes legible and enforceable}

The 2026-07-28 revision is the largest change to the protocol since its launch,
and for a governance argument its importance is specific. After this revision,
every agent action is at once a well-formed record and a policy decision point.
Five properties of the release, taken together, produce that result. None of
them is invented here; each is a mechanism the release already ships, and the
point is what they add up to.

\paragraph{Identity-bound.} Every request now carries the client's identity in
its \texttt{\_meta} block rather than establishing it once per session: the
\texttt{io.modelcontextprotocol/clientInfo} field travels with the call itself
(SEP-2575) \citep{mcpchangelog}. On the authorization side, the release hardens
OAuth so a client must validate a present issuer (\texttt{iss}) parameter against
the recorded authorization server before it redeems a code, closing a class of
mix-up attack (SEP-2468) \citep{mcpchangelog}. An action arrives attributable to
an authenticated principal and a named client implementation.

\paragraph{Typed.} Tool inputs and outputs are lifted to full JSON Schema
2020-12 (SEP-2106), so the shape of every call and its result is declared and
checkable rather than conventional \citep{mcpchangelog}. A record whose payloads
validate against a schema is a record a downstream system can parse without
guessing at its structure.

\paragraph{Traceable.} The release documents W3C Trace Context conventions ---
\texttt{traceparent}, \texttt{tracestate}, and \texttt{baggage} in
\texttt{\_meta} (SEP-414) --- so a single action correlates across SDKs,
gateways, and any OpenTelemetry backend \citep{mcpchangelog}. The deprecation of
the old Logging primitive points the same way: diagnostic logging moves to
\texttt{stderr} and OpenTelemetry (SEP-2577) \citep{deprecated}. Observability
belongs to the tracing plane, not to a bespoke protocol message.

\paragraph{Enforceable.} Because the release removes protocol sessions
(SEP-2567), any request can land on any server replica, and authorization is
therefore evaluated on every call rather than inherited from a session
established once \citep{sep2567}. The release also requires routing headers,
\texttt{Mcp-Method} and \texttt{Mcp-Name}, on Streamable-HTTP POSTs, and requires
servers to reject a request whose headers and body disagree (SEP-2243)
\citep{mcpchangelog}. A gateway can now make an allow, deny, or downgrade
decision on the operation without parsing the body. The per-call boundary is a
real enforcement point.

\paragraph{Fleet-scalable.} List and read results now carry required cache
metadata: a \texttt{ttlMs} freshness hint and a \texttt{cacheScope} of
\texttt{public} or \texttt{private} (SEP-2549) \citep{mcpchangelog}. With
sessions gone, these results are stable enough to cache, which is what turns
governing a fleet of agents from a per-agent tax into a per-server cost.

Set these five properties side by side and the shape is clear. An organization
running agents on the 2026-07-28 spec gets, without extra work, a stream of
actions that are attributable, typed, traceable, individually enforceable, and
cheap to observe at fleet scale. This is a genuine advance in what a protocol
makes governable, and the temptation is to stop here and call the governance
problem solved. It is not solved. A well-formed record has a limit, and that
limit is the subject of the next section.

\section{What the record cannot tell you}

A record can be complete and still be silent about the thing that matters.
Consider an agent action that is fully authorized, correctly typed, cleanly
traced, and schema-valid, a record with nothing wrong in it, that opens a pull
request introducing a SQL injection. Nothing in the protocol layer is equipped to
notice, because nothing in the protocol layer looks at what the change does once
it is merged. Veracode's 2025 analysis found that 45\% of AI-generated code
samples introduced an OWASP Top-10 vulnerability, across more than a hundred
models, eighty tasks, and four languages, with the rate roughly flat regardless
of model size or recency \citep{veracode2025}. The defective change and the
correct change produce identical protocol records.

This is the clean-record trap, and the audit layer has always had it. A log
tells you an action happened and was permitted; it never tells you the action
was any good. The argument of this paper is that conditioning autonomy on task
accuracy re-creates that same trap one level up. Grade an agent on how well it
performs its task, and a workflow that scores 99\% on tool-call correctness still
ships the Veracode defect, because tool-call correctness and delivered-code
quality are different measurements: the defective call \emph{is} a correct call. Worse, an autonomy ladder
built on task accuracy promotes a workflow precisely when its accuracy scores are
high, which is the moment its output is drawing the least human scrutiny. The
metric rises as the examination falls.

Even a perfectly accurate agent would be graded on the wrong variable, because
task accuracy measures the agent's behavior and says nothing about the outcome
that behavior delivered. An industry benchmark of 2.7 million pull requests
across 253 organizations reaches the same conclusion from the field: the output
that matters is delivered work, not generated code, and agent-opened pull
requests merge at markedly lower rates than human-opened ones, because a change
no human owns tends to sit unmerged \citep{aibench2026}. The variable that
counts lives one step past the tool boundary, at the delivered change,
and no protocol call runs there.

\section{The delivery evidence layer}

Start with a limit that most engineering organizations already know they have.
The delivery-metrics programs leaders run today, DORA's four keys and DX Core 4,
measure delivery at the team level, but they cannot tell you whether a movement in
those metrics was driven by AI tooling, by process change, or by a quality
trade-off. DX names this the attribution gap directly: standard delivery metrics,
designed before agents wrote code, cannot say what drove the change they measure
\citep{dxcore4}. An organization can watch its change-failure rate move and have
no way to attribute the move to its agents. The 95\% with no measurable P\&L return and the attribution gap are the
same problem seen from two ends: an outcome that cannot be attributed to the work
that produced it.

Closing that gap requires a layer the protocol does not provide, and naming it
correctly matters, because an adjacent discipline already claimed a
similar-sounding term for the opposite job. Context engineering builds a
\emph{context layer} whose purpose is to feed better inputs to agents so they act
well \citep{ibmcontext}. The layer this paper describes runs in the other
direction. Call it the \emph{delivery evidence layer}: its purpose is to feed
outcome evidence to the authorizer (the human or process that makes the trust
decision, later the cadence) so the organization can decide what an agent
is allowed to do next. Context feeds the actor. Evidence feeds the authorizer.
It is one store, and it has four capabilities.

\paragraph{Attribute.} Attribute joins the spec's per-call action records to
version-control and CI/CD events, so a delivered change carries a legible answer
to which autonomous work produced it and through what path. The join is the
implementer's to build. The spec emits attributable, traceable action records
(SEP-414, SEP-2575) and the pipeline emits commits, reviews, and deploys, but
nothing in the protocol stitches the two together. That stitch is the fix for the
attribution gap, and it is the foundational capability the other three depend on.

\paragraph{Evaluate.} Evaluate attaches outcome signals to each agent-authored
change: change-failure rate, revert, rework, override-before-merge, whether the
change held in production, and cost-to-outcome. These are delivery signals, not
activity counts. They answer a different question than how many actions the agent
took. They answer what the actions delivered, and what the delivery cost.
Evaluate is where the conditioning variable of the whole argument, delivered
outcome, gets computed.

\paragraph{Serve.} Serve hands that evidence back as deterministic, typed,
freshness-contracted context, delivered over the spec's own substrate. Here the
release pays a second dividend: full JSON Schema typing (SEP-2106) makes the
served context checkable, required cache scopes and TTLs (SEP-2549) make its
freshness a contract rather than a hope, and the per-call boundary (SEP-2567) is
where it is delivered and re-authorized. The same protocol that emitted the
events carries the answers back, and Serve reaches two consumers over those rails:
the review cadence (defined in the next section) reads the evidence as the
authorizer's input, and the agents
read the distillation as operating context, the current baselines and constraints
they should act within.

\paragraph{Steer.} Steer feeds evaluations into boundary policy: the outcomes a
workflow has delivered set the autonomy it is granted on its next call.
Steer is the capability that turns the join into a loop, and it is where the
protocol earns its place in this argument, because the per-call enforcement point
is the actuator the loop writes to. It is also the capability most easily
overstated, so the next section defines it carefully and states its limits first.

Attribute and Evaluate are the write path; Serve and Steer are the read path.
None of the four is speculative: a join over data an organization already has,
metrics teams already track, a typed read over a governed transport, and policy
at a boundary the spec now provides. What makes them a layer rather than four
disconnected scripts is that they share one store of delivered-outcome evidence,
keyed to the unit an organization wants to govern.

\section{Earned autonomy}

The 2026-07-28 release ships into a conformance-gated, SDK-tiered ecosystem: a
conformance suite live since 2026-01-23, published SDK tiers since 2026-02-23,
and, as of SEP-2484 (Final 2026-06-04), a standing rule that future
Standards-Track SEPs cannot reach Final without a matching conformance scenario
\citep{sep2484, sdktiers}. The SDK tiers are the part worth studying. A Tier 1
SDK must pass 100\% of the conformance suite; if any test fails continuously for
four weeks, it is relegated to Tier 2; and features move through a lifecycle with
a minimum twelve-month deprecation window measured from the revision release
\citep{sdktiers, mcplifecycle}. Standing is earned by demonstrated conformance
and lost publicly when the demonstration lapses.

That is a working model of earned standing, and it is a model humans run. No
automated controller relegates an SDK. Maintainers publish tiers against measured
conformance, on defined clocks, with a defined relegation rule, and the
governance works. This is the precise shape to port. Earned autonomy is a rolling
trust review in which workflow configurations hold published autonomy tiers,
delivery-evidence-layer signals are the input, promotion and relegation follow
defined triggers, and the resulting grant is enforced per call at the boundary
the spec provides. Steer is that cadence. It is not a controller that reads a
metric and rewrites policy in a closed loop; it is a review rhythm with a clock
and a rule, of the same kind the standard already runs on itself.

Every component of this cadence is mechanically expressible today. Identity and
per-call enforcement come from the spec. Attribution, outcome evaluation, and
served evidence come from the delivery evidence layer. Published tiers, clocks,
and triggers are process. The only genuinely new input is the swap at the center:
replace static scope grants, issued once and rarely revisited, with grants
conditioned on delivered-outcome evidence. Automating the cadence, turning the
rolling human review into a closed control loop, is the forward agenda, not a
claim about today.

Stated plainly, and stated as the sentence a skeptic should attack: autonomy
conditioned on task accuracy re-creates at the trust layer the same clean-record
trap the audit layer already has; delivered outcome is the only conditioning
variable that closes the loop, and the 2026-07-28 spec is the first standard
substrate on which that condition is enforceable per action.

\paragraph{The failure modes are dials, not objections.} Because Steer is a
cadence, its known failure modes are design parameters that set the cadence's
dials. \emph{Attribution lag} --- outcomes arrive days or weeks after the change
--- sets the length of the trailing window the review scores on; instantaneous
scoring is the mistake the window length corrects. \emph{Unstable identity} ---
agent instances are ephemeral and interchangeable --- sets the unit of trust: the
durable thing to govern is a workflow configuration, not an instance, because the
configuration is what produced the work and what will produce it again.
\emph{Goodhart's law} sets the gating signal: gate on held-in-production outcomes
and never on activity counts, because the moment the metric is a throughput count
it is gamed. \emph{Small-n} --- too few observations to trust --- sets the
cold-start default: absent evidence, a workflow starts in a sandbox tier, since
autonomy is earned from evidence and the safe default without it is low.

\paragraph{Two clarifications keep the framework honest} about what the protocol
does and does not provide. First, the per-call boundary is an enforcement point,
not a scope engine. The spec gives the point at which authorization is checked on
every call; it does not decide what a given principal may do. The scope
semantics, the policy itself, remain the implementer's design. Authentication is
not authorization scope. Second, per-call identity attributes to the
authenticated principal and the named client implementation, carried in
\texttt{clientInfo}, and not to an individual agent or subagent. Agent-level
attribution is a convention an organization builds on top, and it is exactly why
the unit of trust is a workflow configuration: that is the granularity the
evidence layer can attribute to, not a per-instance identity the protocol never
reported.

This shape has precedent outside software delivery. Open-source foundations
already promote projects by earned standing: the CNCF moves a project from
Sandbox to Incubating to Graduated on demonstrated adoption, security review, and
governance maturity \citep{cncflifecycle}. The criteria are partly quantitative
and partly qualitative, and the decision is ratified by humans. It is
evidence-based maturation with human ratification, not a purely metric-gated
switch, which is exactly the posture earned autonomy takes toward a workflow
rather than a whole project.

There is also direct empirical support for building the feedback capability the
framework depends on. DORA's 2025 research found that AI acts as an amplifier: it
has a positive relationship with delivery throughput and a negative one with
stability, and which effect dominates is conditional on a team's existing
capability \citep{dora2025}. The organizational property that decides whether AI
speed helps or hurts is the delivery-feedback capability, which is the same
capability this paper argues an organization must build. The amplifier result is
not only a stake. It is corroboration that the conditioning variable is the right
one.

The cadence is configuration, not research. The trust decision --- read the
delivered-outcome evidence, set the tier --- is the organization's to make; the
gateway's only job is to enforce the resulting scope on every call. A workflow
holds the identity its current tier has earned, and the gateway gates each tool
call against that identity:

\begin{figure}[t]
\centering
\begin{minipage}{0.92\linewidth}
\begin{lstlisting}[style=policy]
# Per-call authorization at an MCP gateway. A workflow holds the
# identity its current tier earns; the gateway enforces scope on
# every tool call. Relegation = the cadence reassigns the workflow
# to a lower-tier identity, and its next call is gated anew.
authorization:
  - allow: workflow == "wf-alpha"          # trusted: full access
  - allow: workflow == "wf-beta"           # sandbox: read-only
        && tool in ["read_file", "list_directory", "search_files"]
\end{lstlisting}
\end{minipage}
\caption{\textbf{Validated on a running MCP gateway.} A per-call authorization
rule gating tool scope on a workflow's tier-earned identity. On an isolated
gateway instance, the identical \texttt{write\_file} call was allowed for the
trusted identity and denied per call for the sandbox identity, while the read
call was allowed for both (see \texttt{data/}). The evidence-to-tier decision is
the organization's cadence; the gateway is only the actuator. The trailing window
the cadence scores on absorbs attribution lag, and the unit of trust is a workflow
configuration, not an instance. This is the shape teams already use to gate a
deploy, pointed at delivered-outcome evidence.}
\end{figure}

\section{Recommendations}

The move a leader can make on Monday costs no procurement. It is to stand up the
cadence as a process before it is a product. Concretely: publish autonomy tiers
for the workflow configurations already running agents, from a sandbox tier up
through a trusted tier; define the trailing window and the outcome signals each
tier is scored on; set the relegation triggers and the clock, borrowing the
standard's own four-week shape as a starting point; and default every new or
sparsely-observed workflow to the sandbox. This is a review rhythm a platform or
developer-experience team can run in a spreadsheet on the first day and encode in
gateway policy as it matures. None of it requires buying anything.

The cadence adjusts trust over time; it presupposes a gate that screens every
change regardless of the workflow's record, and that gate is the cheapest control
to stand up. Turn on pre-merge security scanning, blast-radius and impact gating,
and agentic code review on every pull request, agent-authored or not. Agentic code
review is the fastest of the three to roll out org-wide and the lowest-friction:
the same industry benchmark found that adding it lifted the share of pull requests
merged within thirty days by up to about five percentage points, and that it
recovers much of the yield agents lose when they open changes no human owns
\citep{aibench2026}. This is the preventive layer earned autonomy sits on; it is
not an alternative to the cadence.

The capability to build alongside it is the delivery evidence layer, and it
belongs in the same budget cycle as the MCP boundary work, not after it. The
boundary makes actions enforceable; the layer makes outcomes legible; the cadence
needs both. Sequencing the enforcement work first and the evidence work later
produces the exact state this paper opened on: a fleet whose actions are
perfectly recorded and whose value is unattributable.

When evaluating any system, internal or purchased, that claims to govern agent
autonomy, five questions separate an outcome-conditioned tool from an
activity-graded one:

\begin{itemize}[leftmargin=1.4em,itemsep=0.25em,topsep=0.3em]
  \item \textbf{Attribute at the change level.} Does it join governed actions to
  the commit, review, and deploy they produced, or does it stop at the action
  record?
  \item \textbf{Close the governed-action-to-shipped-result loop.} Can a granted
  scope be traced to the outcome it caused, so authority and result meet in one
  place?
  \item \textbf{Segment authorship.} Does it tell agent-authored change apart
  from human-authored change, so the attribution gap is closed rather
  than aggregated over?
  \item \textbf{Resist adoption proxies.} Does it condition trust on
  held-in-production outcomes and refuse activity and adoption counts, such as
  action volume or acceptance rate, as stand-ins for delivery?
  \item \textbf{Serve context over governed rails.} Does it return evidence to
  the agents over the same typed, cache-contracted, per-call-authorized transport
  the actions ran on?
\end{itemize}

The 2026-07-28 spec settles the question of whether an organization can see and
control what its agents do, action by action. It leaves open the question this
paper is about: whether the organization can tell what those actions delivered,
and condition trust on the answer. The substrate for that second question now
exists. The layer above it, and the cadence above that, are the work in front of
every team running agents at scale.


\bibliographystyle{plainnat}
\bibliography{references}

\end{document}
